% TOP_PROBLEM: {"id": 3000067, "title": "Quasi-kernels and quasi-sinks", "category_id": 2, "category": {"id": 2, "name": "combinatorics", "display_name": "Combinatorics", "description": "Counting problems, graph theory, discrete structures.", "slug": "combinatorics", "order_index": 2, "created_at": "2026-07-31T15:26:25.671Z"}, "status": "open"}
% FIRST_POSTED: null
% ENTRY_KIND: statement_only
% Imported from: batch09.tex; UnsolvedMath ID 3000067
\documentclass[11pt]{article}
\usepackage[margin=25mm]{geometry}
\usepackage{fontspec}
\setmainfont{Latin Modern Roman}
\usepackage{amsmath,amssymb,mathtools}
\usepackage{unicode-math}
\setmathfont{Latin Modern Math}
\usepackage{xurl,hyperref}
\hypersetup{colorlinks=true,urlcolor=blue}
\setcounter{secnumdepth}{0}
\setlength{\parindent}{0pt}
\setlength{\parskip}{5pt}
\setlength{\emergencystretch}{3em}
\newcommand{\sref}[2]{\hyperlink{source:#1:#2}{[S#2]}}
\newcommand{\cataloguescope}[1]{\paragraph{Recorded target.}#1}
\begin{document}
% UnsolvedMath ID 3000067; source catalogue code AMR-029-0067
\setcounter{section}{3000066}
\section{Quasi-kernels and quasi-sinks}\label{3000067}
{\small\sffamily UnsolvedMath ID 3000067 · Combinatorics}
\subsection{Definitions and mathematical statement}

\paragraph{Shared notation.}$\mathbb N=\{1,2,\ldots\}$, and $\mathbb Z,\mathbb Q,\mathbb R,\mathbb C$ denote the integers, rationals, reals and complex numbers. Any other conventions are specified locally.


Let $D=(V,A)$ be a digraph, where $V$ is any infinite set. A set is independent if no arc joins two of its vertices. In this question a quasi-kernel $K$ is an independent set such that every vertex is reachable from some member of $K$ by a finite directed path of length at most two. A quasi-sink is a quasi-kernel after all arcs are reversed. Paths of length zero are allowed. For $W\subseteq V$, $D[W]$ denotes the induced digraph.
\paragraph{Statement recorded in the supplied source.}
Is there always a partition $V=V_1\dot\cup V_2$ such that $D[V_1]$ has a quasi-kernel and $D[V_2]$ has a quasi-sink? Either part may be empty, in which case its empty set supplies the required object. \sref{3000067}{2}
\subsection{Short English statement}
Can every infinite digraph be partitioned into an induced part dominated in two steps by an independent set and a part satisfying the reversed condition?
\subsection{Sources}
\begin{description}
\item[\hypertarget{source:3000067:1}{S1}] ULAM.AI and the UnsolvedMath Contributors, \emph{UnsolvedMath: A Curated Collection of Open Mathematics Problems}, record 3000067 (AMR-029-0067). CC BY 4.0. \url{https://huggingface.co/datasets/ulamai/UnsolvedMath}.
\item[\hypertarget{source:3000067:2}{S2}] P. L. Erdős and L. Soukup, Quasi-kernels and quasi-sinks in infinite graphs (2007), Introduction, definitions and Conjecture 1.2, pp.1--2. \url{https://arxiv.org/pdf/0712.0663}.
\end{description}
\label{end:3000067}
\end{document}
